\honest{Angry Atoms}{Eliezer Yudkowsky}{2008}

No one can see how quarks make up a hand; at best you know the experiments, at worst you were told. But hands push on things, and we picture atoms as little billiard balls that push on things too, so a hand made of atoms is easy to imagine.

But what about anger? ``How could little billiard balls be angry? Tiny frowny faces on the billiard balls?'' \nb{The ancient atomists tried. Lucretius gave the mind an element of heat and wrote that lions have ``more of hot'' than stags, which explains anger by a kind of particle.} Put yourself in the place of a hunter-gatherer with no idea of computation or neurons. Set experience aside and look only at behaviour: anger makes people hit, gossip and plant scorpions in tents, while billiard balls only push. The gap looks uncrossable. \nb{The phrase ``explanatory gap'' is Joseph Levine's, coined in 1983 for the gap between physical facts and how an experience feels, the case the post sets aside.}

Then I play the objector: how would the billiard balls know how to plot, unless they were angry themselves? ``Maybe wine just contains a potency of angerness. Clearly, reductionism is just a flawed notion.''

Saying that neurons ``process information'' does not cross the gap. What crosses it is knowing how to build a chess program: search the tree of moves, score the positions at the end, and trace back with minimax. More generally, a mind whose inner causal chains mirror the world can run a utility function over what it imagines and choose an action. \nb{Kenneth Craik described minds that try out alternatives on a ``small-scale model'' of reality in 1943, and Claude Shannon proposed minimax search for chess in 1950. The post does not claim the idea as new.} ``All this is still tremendously oversimplified,'' I admit, but now you can see how alcohol makes a planner angrier: ``The billiard balls in the alcohol push on the billiard balls making up the utility function.'' \nb{That sentence is the whole account of what makes a plan an angry one.}

No one can visualize the step from neurons to anger, as one can a hand made of fingers. Someone who only repeats ``Anger is hormones'' ends up saying that anger is unjustified, or that it does not exist. ``I think this is what non-reductionists/non-materialists think they are criticizing.'' \nb{No one is quoted. David Chalmers, a leading non-materialist whose zombie argument the book takes up later, called explaining how a system controls behaviour an ``easy problem'', with ``no obvious obstacle'' to a computational explanation. Chalmers's objection was about experience.}

This was an easy example, I grant, because it concerns behaviour and not experience: ``just a practice problem''. Still, ``Explanatory gaps can be crossed, if you accept help from science''.
